\documentclass[14pt,a4paper]{extarticle}
\usepackage{lipsum} % can be remove
\usepackage{geometry}
 \geometry{
 a4paper,
 total={170mm,257mm},
 left=2.54cm,
 right=2.54cm,
 top=2.54cm,
 bottom=2.54cm,
 }
%--ตั้งหัวกระดาษ--%
\usepackage{fancyhdr}
\pagestyle{fancyplain}

\fancyhead[L]{การทดลองหาประสิทธิภาพด้านเวลาและการป้องกันการโจมตีของ RYU SDN Framework}
\fancyhead[R]{\thepage}
%--*********--%
\usepackage{comment}
\usepackage[no-math]{fontspec} % จัดการเกี่ยวกับฟ้อนต์
\usepackage{graphicx} % จัดการเกี่ยวกับภาพ
\graphicspath{ {./img/} }
%--สำหรับภาษาไทยและการตัดบรรทัด--%
\XeTeXlinebreaklocale "th_TH"
\XeTeXlinebreakskip = 0pt plus 1pt  
\setmainfont[
  Path=fonts/,
  BoldFont={THSarabunNew_Bold.ttf},
  ItalicFont={THSarabunNew_Italic.ttf},
  BoldItalicFont={THSarabunNew_BoldItalic.ttf},
]{THSarabunNew.ttf}
%--*********************--%
%--config--%
\usepackage{pgfplots} 
\usepackage{glossaries}
\newacronym{sdn}{ระบบเครือข่ายที่กําหนดโดยซอฟต์แวร์}{ Software-defined networking }
\newacronym{ryu}{RYU}{ RYU SDN framework }
\newacronym{software}{ซอฟต์แวร์}{ Software }
\newacronym{network-system}{ระบบเครือข่าย}{ Network system }
\usepackage{titlesec}
\titleformat{\section}
  {\normalfont\fontsize{16}{15}\bfseries}{\thesection}{1em}{}
\titleformat{\subsection}
  {\normalfont\fontsize{16}{15}\bfseries}{\thesubsection}{1em}{}
\renewcommand{\figurename}{รูปภาพที่}
%
\titlespacing*{\section}
{0pt}{5.5ex plus 1ex minus .2ex}{0ex plus 0ex}
\titlespacing*{\subsection}
{0pt}{0ex}{0ex}

%
\begin{document}
%--สำหรับหน้าปก--%
\begin{titlepage}
  \begin{figure}
    \centering
    \includegraphics[width=1.5cm]{Ubu_logo.png}
  \end{figure}
  \begin{center}
    \LARGE การทดลองหาประสิทธิภาพด้านเวลาและการป้องกันการโจมตีของ RYU SDN Framework โดยใช้ Group table และ Proxy arp ในรูปแบบ Single Controller และ Multi Controllers\\
    \vspace{6cm} 
    \LARGE ณวสันต์  วิศิษฏ์ศิงขร \\
    \vfill
    \LARGE แขนงวิชาเทคโนโลยีเครือข่ายและความมั่นคงทางไซเบอร์ \\ 
    สาขาวิชาเทคโนโลยีสารสนเทศและการสื่อสาร \\
    คณะวิทยาศาสตร์
    มหาวิทยาลัยอุบลราชธานี
    \pagebreak
  \end{center}
\end{titlepage}
%--***************--%
% \section{ระเบียบวิธี} 

\subsection{วัตถุประสงค์}  
\gls{sdn} นั้นมีการใช้งานและการตั้งค่าที่หลากหลาย
ไม่ว่าจะเพื่อให้เหมาะกับ \gls{network-system} ขององค์กรหรือหน่วยงานของตน ยังต้องช่วยเพิ่มประสิทธิภาพของ\gls{network-system}อีกด้วย
นอกจากนั้น \gls{sdn} ยังถูกนำมาใช้เพื่อป้องกันการโจมตีผ่าน\gls{network-system} ไม่ว่าจะเป็น Dos หรือ DDos
อย่างไรก็ตาม\gls{sdn}นั้นก็คือ \gls{software} รูปแบบหนึ่ง
ดังนั้นก็จะมี\gls{software}หลายตัวที่ถูกพัฒนาขึ้นเป็น ``\gls{sdn}'' เช่น RYU OpenDaylight NOX POX ฯลฯ 
ดังนั้น
วัตถุประสงค์ของการทดลองครั้งนี้ก็เพื่อหาประสิทธิภาพทั้งในด้านเวลาและการตรวจจับการโจมตีของ \gls{ryu} ซึ่งเป็น\gls{sdn}ประเภทหนึ่ง ที่ถูกรันอยู่บน SDN Controller
พร้อมกับการใช้งาน Group table หรือ Proxy ARP ทั้งในรูปแบบ Single Controller และ Multi Controllers
\\
\subsection{การตั้งค่าการจำลอง (Simulation setup)}

\textbf{เครื่องมือและอุปกร์ที่ใช้ในการทดลองมีดังนี้} 

\notebook รุ่น Lenovo IdeaPad 5 14IIL05 รันด้วยระบบปฏิบัติการ Arch Linux 6.7.4-arch1-1 
เป็นอุปกรณ์ที่ใช้สำหรับการรัน\gls{software}ทดสอบทั้งหมด

GNOME Terminal Version 3.50.1 
เป็น\gls{software}อยู่บน\notebook ที่จะใช้ในการ ssh เข้าไปใน server เพื่อทำการสั่งการทำงานต่างๆ

Oracle VM VirtualBox เวอร์ชั่น 7.0.14 
ใช้สำหรับการจำลอง \gls{server} ขึ้นมาเป็นอีกเครื่องหนึ่งบน\notebook
ที่กล่าวไปข้างต้น ซึ่งจะทำให้เสมือนว่ามี\gls{server}เพิ่มขึ้นมาอีกเครื่องหนึ่ง

ระบบปฏิบัติการ Ubuntu server 22.04.3
คือระบบปฏิบัติการที่ถูกรันอยู่บน\gls{server}ที่จำลองอยู่บน Oracle VM VirtualBox 

Python version 3.9.18 และ 2.7.18 
คือรันไทม์สำหรับรันสคริปต์ไฟล์ที่เขียนด้วยภาษา Python ซึ่งจะใช้ในการรัน \gls{ryu} รวมไปถึงการจำลอง Topology ของ mininet

RYU version 4.34
คือ \gls{sdn} ที่ได้กล่าวไปข้างต้น ซึ่งเป็น\gls{software}หลักตัวหนึ่ง ที่จะถูกใช้ในการรันเชื่อมต่อกับ \gls{controller} บน mininet

mininet version 2.3.1b4
คือ\gls{software}ที่จะทำการจำลอง Topology ที่ใช้ในการทดสอบขึ้นมาภายในเครื่องจำลอง Ubuntu server ที่รันอยู่บน Oracle VM VirtualBox

wireshark 3.6.2
คือ\gls{software}ที่จะช่วยในการดักจับข้อมูลที่วิ่งผ่าน\gls{network-system} ซึ่งจะเป็น\gls{software}ที่ใช้ในการวัดประสิทธิภาพในด้านเวลาและการป้องกันการโจมตี

vim version 8.2.2121
คือ\gls{software}ที่ใช้ในการแก้ไขไฟล์ข้อความหรือไฟล์ข้อความต่างๆ บนระบบปฏิบัติการ Ubuntu server

Vscode 1.86.1
เป็น\gls{software}อีกตัวหนึ่งที่ใช้ในการแก้ไขไฟล์ข้อความซึ่งมีส่วนเสริมที่ช่วยในการเขียน ซึ่งในการทดลองนี้จะนำมาใช้ในการแก้ไขไฟล์สคริปต์ของ Python ในกรณีที่ต้องการความแม่นยำหรือแก้ปัญหาที่ซับซ้อน
\newline
\subsection{รูปแบบในการทดลอง}
ในการจะทำการทดลองนั้นจำเป็นต้องมี python script ที่ใช้ในการสร้าง topology จำลองบน mininet
ซึ่งต้องสร้างทั้งแบบ Single Controller และ Multi Controllers โดยสามารถใช้ mininet gui สร้างได้ ตามภาพด้านล่างนี้
หรือสามารถดาวน์โหลดไฟล์ scripts ได้ที่ http://projectcs.sci.ubu.ac.th/nawasan/sdn-topo-mininet
\\

\begin{figure}[h!]
    \centering
    \includegraphics*[width=20em]{single-edit.png}
    \caption{Topology แบบ Single Controller}
    \label{img:topo_cs}
\end{figure}

\begin{figure}[h!]
    \centering
    \includegraphics*[width=20em]{multi-edit.png}
    \caption{Topology แบบ Multi Controllers}
    \label{img:topo_cm}
\end{figure}

\pagebreak

โดยรูปภาพที่ \ref{img:topo_cs} คือ Topology แบบมี 1 \gls{controller} โดยจะทำการกำหนดให้\gls{controller}เป็นแบบ remote 
และรูปภาพที่ \ref{img:topo_cm} คือ Topology แบบมีหลาย\gls{controller} โดยจะทำการกำหนดให้\gls{controller}เป็นแบบ remote ทั้งหมด และกำหนด port เป็น 6633 6634 และ 6635
นอกจากนั้น เมื่อสร้าง script python ขึ้นมาแล้ว หรือดาวน์โหลดไฟล์จากลิงก์ที่ให้ไว้ จะทำการตั้งชื่อไฟล์เป็นดังนี้ Topology ที่มี 1 \gls{controller}
ดังรูปภาพที่ \ref{img:topo_cs} จะทำการตั้งชื่อเป็น single-topo.py และ Topology ที่มีหลาย\gls{controller} ดังรูปภาพที่ \ref{img:topo_cm} จะทำการตั้งชื่อเป็น multi-topo.py

ในการทดลองครั้งนี้ผู้ทำการทดลองได้ออกแบบการทดลองให้มีความเข้มงวดและรัดกุมมากที่สุด เพื่อให้ผลลัพธ์ที่ได้มีความคลาดเคลื่อนน้อยที่สุด
โดยผู้ทำการทดลองจะออกแบบการทดลองโดยให้มี 2 รูปแบบ 4 เงื่อนไข ดังนี้


เงื่อนไขที่ 1. (A1) group table แบบ 1 \gls{controller}\\
โดยจะทำการรัน ryu-manager ด้วย script ไฟล์ที่ชื่อ ex7\_group\_tables.py 
และทำการรันไฟล์ชื่อ single-topo.py ด้วย python เพื่อจำลอง Topology แบบ 1 \gls{controller}


เงื่อนไขที่ 2. (A2) group table แบบหลาย\gls{controller} \\
โดยจะทำการรัน ryu-manager ด้วย script ไฟล์ที่ชื่อ ex7\_group\_tables.py 
และทำการรันไฟล์ชื่อ multi-topo.py ด้วย python เพื่อจำลอง Topology แบบหลาย\gls{controller}


เงื่อนไขที่ 3. (A3) proxy arp แบบ 1 \gls{controller}\\
โดยจะทำการรัน ryu-manager ด้วย script ไฟล์ที่ชื่อ ex8\_arp\_proxy.py
และทำการรันไฟล์ชื่อ single-topo.py ด้วย python เพื่อจำลอง Topology แบบ 1 \gls{controller}

เงื่อนไขที่ 4. (A4) proxy arp แบบหลาย\gls{controller} \\
โดยจะทำการรัน ryu-manager ด้วย script ไฟล์ที่ชื่อ ex8\_arp\_proxy.py 
และทำการรันไฟล์ชื่อ multi-topo.py ด้วย python เพื่อจำลอง Topology แบบหลาย\gls{controller}

รูปแบบที่ 1. คือ ทำการทดสอบโดยจับแพ็กเก็ตทีละจำนวน 100,000 250,000 500,000 750,000 และ 1,000,000 แพ็กเก็ต
จากนั้นนำมาวิเคราะห์เปรียบเทียบหาประสิทธิภาพด้านเวลาและการตรวจจับการโจมตี

รูปแบบที่ 2. คือ ทำการทดสอบโดยจำแพ็กเก็ตจำนวน 1,000,000 และ 2,000,000 แพ็กเก็ต
จากนั้นนำมาวิเคราะห์เปรียบเทียบหาประสิทธิภาพด้านเวลาและการตรวจจับการโจมตี
\\

สาเหตุที่มีการทดลองสองรูปแบบนั้นก็เพื่อให้สามารถตรวจสอบได้ว่าผลลัพธ์ที่ได้จากการทดสอบมีความถูกต้องน่าเชื่อถือ
เป็นต้นว่าหากผลการทดสองจากทั้งสองรูปแบบสอดคล้องกัน ก็สามารถเชื่อได้ว่าผลการทดลองนั้นถูกต้อง นอกจากนั้นผู้ทำการทดลองยังได้กำหนดให้ลำดับวิธีของทั้ง 2 รูปแบบการทดลองมีความต่างกันในบางขั้นตอน เพื่อเพิ่มความเชื่อมั่นและความถูกต้องมากขั้นไปอีก
\\
\subsection{ขั้นตอนลำดับวิธีทดลอง}
ในการทดลอง ขั้นตอนลำดับจะต่างกันในส่วนของรูปแบบการทดลองเท่านั้น โดยที่เงื่อนไขต่างๆ จะยังคงมีลำดับการทดลองที่เหมือนกัน กล่าวคือในเงื่อนไขการทดลองนั้นจะต่างกันก็เพียงแต่ไฟล์ที่ทำการรันเท่านั้น โดยที่ลำดับขั้นตอนจะยังคงเหมือนเดิม 
ในแต่ละเงิ่อนไขการทดลองจะทำการรันทดสอบ 3 รอบ และเลือกค่ากลางมาทำการวิเคราะห์ กล่าวคือ เลือกค่าที่ไม่น้อยที่สุดและไม่มากที่สุดจากค่าที่ได้จากการทดสอบ

\subsubsection*{รูปแบบการทดลองที่ 1 บน 1 \gls{controller}}

ในขั้นต้นผู้ทำการทดลองต้องทำการเปิด gnome terminal ขึ้นมาและเปิดหน้าต่างย่อยเป็นจำนวน 3 หน้าต่าง
ทุกหน้าต่างจะต้องทำการเชื่อมต่อไปยัง\gls{server}ด้วย ssh ให้เรียบร้อย
โดย หน้าต่างที่ 1 จะทำการรัน ryu-manager หน้าต่างที่ 2 จะทำการรัน mininet และหน้าต่างที่ 3 จะทำการรัน wireshark
โดยมีลำดับขั้นตอนการทดลองดังนี้


1. หน้าต่างที่ 1 ทำการรันคำสั่ง    
\$ ryu-manager learn-sdn-with-ryu/ryu-exercises/<python script> \\
โดย <python script> จะแทนด้วย ex7\_group\_tables.py หรือ ex8\_arp\_proxy.py ตามเงื่อนไขการทดลอง

2. หน่าต่างที่ 2 ทำการรันคำสั่ง    
\$ sudo python3 single\_topo.py \\
เพื่อทำการจำลอง topology แบบ 1 \gls{controller} ขึ้นมา

3. หน้าต่างที่ 3 ทำการรันคำสั่ง 
\$ sudo wireshark \\
โดยทำการจับแพ็กเก็ตที่ s4-eth1 (ที่เชื่อมกับ h2) และ s11-eth2 (ที่เชื่อมกับ h12) เป็นจำนวนที่เคยกล่าวไปข้างต้น

4. หลังจากที่ wireshark ทำการเริ่มจับแพ็กเก็ต หน้าต่างที่ 2 จะทำการรันคำสั่ง \$ mininet> h2 hping3 h12 -S --flood -V
โดยทันที เพื่อทำการให้ h2 โจมตีไปที่ h12

5. หลังจากทำการจับแพ็กเก็ตครบตามจำนวนที่กำหนด หน้าต่างที่ 2 ทำการหยุดการโจมตี จากนั้นทำการเก็บค่าสถิติ และทำการปิดโปรแกรม wireshark
หน้าต่างที่ 2 ทำการออกจาก mininet และรันคำสั่ง \$ sudo mn -c หากมีการทดสอบรอบต่อไปต้องเริ่มจากขั้นตอนที่ 1 ใหม่เท่านั้น

\subsubsection*{รูปแบบการทดลองที่ 1 บนหลาย\gls{controller}}

ในขั้นต้นผู้ทำการทดลองต้องทำการเปิด gnome terminal ขึ้นมาและเปิดหน้าต่างย่อยเป็นจำนวน 5 หน้าต่าง
ทุกหน้าต่างจะต้องทำการเชื่อมต่อไปยัง\gls{server}ด้วย ssh ให้เรียบร้อย
โดย หน้าต่างที่ 1 2 และ 3 จะทำการรัน ryu-manager หน้าต่างที่ 4 จะทำการรัน mininet และหน้าต่างที่ 5 จะทำการรัน wireshark
โดยมีลำดับขั้นตอนการทดลองดังนี้

1. หน้าต่างที่ 1 ทำการรันคำสั่ง \\
\$ ryu-manager learn-sdn-with-ryu/ryu-exercises/<python script> --ofp-tcp-listen-port 6633 \\
\indent หน้าต่างที่ 2 ทำการรันคำสั่ง \\
\$ ryu-manager learn-sdn-with-ryu/ryu-exercises/<python script> --ofp-tcp-listen-port 6634  \\
\indent หน้าต่างที่ 3 ทำการรันคำสั่ง \\
\$ ryu-manager learn-sdn-with-ryu/ryu-exercises/<python script> --ofp-tcp-listen-port 6635  \\
โดย <python script> จะแทนด้วย 
ex7\_group\_tables.py หรือ ex8\_arp\_proxy.py
ตามเงื่อนไขการทดลอง

2. หน่าต่างที่ 4 ทำการรันคำสั่ง    
\$ sudo python3 multi\_topo.py \\
เพื่อทำการจำลอง topology แบบหลาย\gls{controller} ขึ้นมา

3. หน้าต่างที่ 5 ทำการรันคำสั่ง 
\$ sudo wireshark 
โดยทำการจับแพ็กเก็ตที่ s4-eth1 (ที่เชื่อมกับ h2) และ s11-eth2 (ที่เชื่อมกับ h12) เป็นจำนวนที่เคยกล่าวไปข้างต้น

4. หลังจากที่ wireshark ทำการเริ่มจับแพ็กเก็ต หน้าต่างที่ 4 จะทำการรันคำสั่ง \$ mininet> h2 hping3 h12 -S --flood -V
โดยทันที เพื่อทำการให้ h2 โจมตีไปที่ h12

5. หลังจากทำการจับแพ็กเก็ตครบตามจำนวนที่กำหนด หน้าต่างที่ 4 ทำการหยุดการโจมตี จากนั้นทำการเก็บค่าสถิติ และทำการปิดโปรแกรม wireshark
หน้าต่างที่ 2 ทำการออกจาก mininet และรันคำสั่ง \$ sudo mn -c หากมีการทดสอบรอบต่อไปต้องเริ่มจากขั้นตอนที่ 1 ใหม่เท่านั้น

\subsubsection*{รูปแบบการทดลองที่ 2 บน 1 \gls{controller}}

ในขั้นต้นผู้ทำการทดลองต้องทำการเปิด gnome terminal ขึ้นมาและเปิดหน้าต่างย่อยเป็นจำนวน 3 หน้าต่าง
ทุกหน้าต่างจะต้องทำการเชื่อมต่อไปยัง\gls{server}ด้วย ssh ให้เรียบร้อย
โดย หน้าต่างที่ 1 จะทำการรัน ryu-manager หน้าต่างที่ 2 จะทำการรัน mininet และหน้าต่างที่ 3 จะทำการรัน wireshark
โดยมีลำดับขั้นตอนการทดลองดังนี้


1. หน้าต่างที่ 1 ทำการรันคำสั่ง    
\$ ryu-manager learn-sdn-with-ryu/ryu-exercises/<python script> \\
โดย <python script> จะแทนด้วย ex7\_group\_tables.py หรือ ex8\_arp\_proxy.py ตามเงื่อนไขการทดลอง

2. หน่าต่างที่ 2 ทำการรันคำสั่ง    
\$ sudo python3 single\_topo.py \\
เพื่อทำการจำลอง topology แบบ 1 \gls{controller} ขึ้นมา

3. หน้าต่างที่ 2 จะทำการรันคำสั่ง \$ mininet> h2 hping3 h12 -S --flood -V
เพื่อทำการให้ h2 โจมตีไปที่ h12

4. หน้าต่างที่ 3 ทำการรันคำสั่ง 
\$ sudo wireshark 
จากนั้นจะทำการเริ่มจับแพ็กเก็ตหลังจากเริ่มการโจมตีแล้ว 30 วินาที
โดยทำการจับแพ็กเก็ตที่ s4-eth1 (ที่เชื่อมกับ h2) และ s11-eth2 (ที่เชื่อมกับ h12) เป็นจำนวนที่เคยกล่าวไปข้างต้น


5. หลังจากทำการจับแพ็กเก็ตครบตามจำนวนที่กำหนด หน้าต่างที่ 2 ทำการหยุดการโจมตี จากนั้นทำการเก็บค่าสถิติ และทำการปิดโปรแกรม wireshark
หน้าต่างที่ 2 ทำการออกจาก mininet และรันคำสั่ง \$ sudo mn -c หากมีการทดสอบรอบต่อไปต้องเริ่มจากขั้นตอนที่ 1 ใหม่เท่านั้น

\subsubsection*{รูปแบบการทดลองที่ 2 บนหลาย\gls{controller}}

ในขั้นต้นผู้ทำการทดลองต้องทำการเปิด gnome terminal ขึ้นมาและเปิดหน้าต่างย่อยเป็นจำนวน 5 หน้าต่าง
ทุกหน้าต่างจะต้องทำการเชื่อมต่อไปยัง\gls{server}ด้วย ssh ให้เรียบร้อย
โดย หน้าต่างที่ 1 2 และ 3 จะทำการรัน ryu-manager หน้าต่างที่ 4 จะทำการรัน mininet และหน้าต่างที่ 5 จะทำการรัน wireshark
โดยมีลำดับขั้นตอนการทดลองดังนี้

1. หน้าต่างที่ 1 ทำการรันคำสั่ง \\
\$ ryu-manager learn-sdn-with-ryu/ryu-exercises/<python script> --ofp-tcp-listen-port 6633 \\
\indent หน้าต่างที่ 2 ทำการรันคำสั่ง \\
\$ ryu-manager learn-sdn-with-ryu/ryu-exercises/<python script> --ofp-tcp-listen-port 6634  \\
\indent หน้าต่างที่ 3 ทำการรันคำสั่ง \\
\$ ryu-manager learn-sdn-with-ryu/ryu-exercises/<python script> --ofp-tcp-listen-port 6635  \\
โดย <python script> จะแทนด้วย 
ex7\_group\_tables.py หรือ ex8\_arp\_proxy.py
ตามเงื่อนไขการทดลอง

2. หน่าต่างที่ 4 ทำการรันคำสั่ง    
\$ sudo python3 multi\_topo.py เพื่อทำการจำลอง topology แบบหลาย\gls{controller} ขึ้นมา

4. หน้าต่างที่ 4 จะทำการรันคำสั่ง \$ mininet> h2 hping3 h12 -S --flood -V
เพื่อทำการให้ h2 โจมตีไปที่ h12

3. หน้าต่างที่ 5 ทำการรันคำสั่ง 
\$ sudo wireshark 
จากนั้นจะทำการเริ่มจับแพ็กเก็ตหลังจากเริ่มการโจมตีแล้ว 30 วินาที
โดยทำการจับแพ็กเก็ตที่ s4-eth1 (ที่เชื่อมกับ h2) และ s11-eth2 (ที่เชื่อมกับ h12) เป็นจำนวนที่เคยกล่าวไปข้างต้น

5. หลังจากทำการจับแพ็กเก็ตครบตามจำนวนที่กำหนด หน้าต่างที่ 4 ทำการหยุดการโจมตี จากนั้นทำการเก็บค่าสถิติ และทำการปิดโปรแกรม wireshark
หน้าต่างที่ 2 ทำการออกจาก mininet และรันคำสั่ง \$ sudo mn -c หากมีการทดสอบรอบต่อไปต้องเริ่มจากขั้นตอนที่ 1 ใหม่เท่านั้น
\\\\

\section{ระเบียบวิธี} 

\subsection{วัตถุประสงค์}  
\gls{sdn} นั้นมีการใช้งานและการตั้งค่าที่หลากหลาย
ไม่ว่าจะเพื่อให้เหมาะกับ \gls{network-system} ขององค์กรหรือหน่วยงานของตน ยังต้องช่วยเพิ่มประสิทธิภาพของ\gls{network-system}อีกด้วย
นอกจากนั้น \gls{sdn} ยังถูกนำมาใช้เพื่อป้องกันการโจมตีผ่าน\gls{network-system} ไม่ว่าจะเป็น Dos หรือ DDos
อย่างไรก็ตาม\gls{sdn}นั้นก็คือ \gls{software} รูปแบบหนึ่ง
ดังนั้นก็จะมี\gls{software}หลายตัวที่ถูกพัฒนาขึ้นเป็น ``\gls{sdn}'' เช่น RYU OpenDaylight NOX POX ฯลฯ 
ดังนั้น
วัตถุประสงค์ของการทดลองครั้งนี้ก็เพื่อหาประสิทธิภาพทั้งในด้านเวลาและการตรวจจับการโจมตีของ \gls{ryu} ซึ่งเป็น\gls{sdn}ประเภทหนึ่ง ที่ถูกรันอยู่บน SDN Controller
พร้อมกับการใช้งาน Group table หรือ Proxy ARP ทั้งในรูปแบบ Single Controller และ Multi Controllers
\\
\subsection{การตั้งค่าการจำลอง (Simulation setup)}

\textbf{เครื่องมือและอุปกร์ที่ใช้ในการทดลองมีดังนี้} 

\notebook รุ่น Lenovo IdeaPad 5 14IIL05 รันด้วยระบบปฏิบัติการ Arch Linux 6.7.4-arch1-1 
เป็นอุปกรณ์ที่ใช้สำหรับการรัน\gls{software}ทดสอบทั้งหมด

GNOME Terminal Version 3.50.1 
เป็น\gls{software}อยู่บน\notebook ที่จะใช้ในการ ssh เข้าไปใน server เพื่อทำการสั่งการทำงานต่างๆ

Oracle VM VirtualBox เวอร์ชั่น 7.0.14 
ใช้สำหรับการจำลอง \gls{server} ขึ้นมาเป็นอีกเครื่องหนึ่งบน\notebook
ที่กล่าวไปข้างต้น ซึ่งจะทำให้เสมือนว่ามี\gls{server}เพิ่มขึ้นมาอีกเครื่องหนึ่ง

ระบบปฏิบัติการ Ubuntu server 22.04.3
คือระบบปฏิบัติการที่ถูกรันอยู่บน\gls{server}ที่จำลองอยู่บน Oracle VM VirtualBox 

Python version 3.9.18 และ 2.7.18 
คือรันไทม์สำหรับรันสคริปต์ไฟล์ที่เขียนด้วยภาษา Python ซึ่งจะใช้ในการรัน \gls{ryu} รวมไปถึงการจำลอง Topology ของ mininet

RYU version 4.34
คือ \gls{sdn} ที่ได้กล่าวไปข้างต้น ซึ่งเป็น\gls{software}หลักตัวหนึ่ง ที่จะถูกใช้ในการรันเชื่อมต่อกับ \gls{controller} บน mininet

mininet version 2.3.1b4
คือ\gls{software}ที่จะทำการจำลอง Topology ที่ใช้ในการทดสอบขึ้นมาภายในเครื่องจำลอง Ubuntu server ที่รันอยู่บน Oracle VM VirtualBox

wireshark 3.6.2
คือ\gls{software}ที่จะช่วยในการดักจับข้อมูลที่วิ่งผ่าน\gls{network-system} ซึ่งจะเป็น\gls{software}ที่ใช้ในการวัดประสิทธิภาพในด้านเวลาและการป้องกันการโจมตี

vim version 8.2.2121
คือ\gls{software}ที่ใช้ในการแก้ไขไฟล์ข้อความหรือไฟล์ข้อความต่างๆ บนระบบปฏิบัติการ Ubuntu server

Vscode 1.86.1
เป็น\gls{software}อีกตัวหนึ่งที่ใช้ในการแก้ไขไฟล์ข้อความซึ่งมีส่วนเสริมที่ช่วยในการเขียน ซึ่งในการทดลองนี้จะนำมาใช้ในการแก้ไขไฟล์สคริปต์ของ Python ในกรณีที่ต้องการความแม่นยำหรือแก้ปัญหาที่ซับซ้อน
\newline
\subsection{รูปแบบในการทดลอง}
ในการจะทำการทดลองนั้นจำเป็นต้องมี python script ที่ใช้ในการสร้าง topology จำลองบน mininet
ซึ่งต้องสร้างทั้งแบบ Single Controller และ Multi Controllers โดยสามารถใช้ mininet gui สร้างได้ ตามภาพด้านล่างนี้
หรือสามารถดาวน์โหลดไฟล์ scripts ได้ที่ http://projectcs.sci.ubu.ac.th/nawasan/sdn-topo-mininet
\\

\begin{figure}[h!]
    \centering
    \includegraphics*[width=20em]{single-edit.png}
    \caption{Topology แบบ Single Controller}
    \label{img:topo_cs}
\end{figure}

\begin{figure}[h!]
    \centering
    \includegraphics*[width=20em]{multi-edit.png}
    \caption{Topology แบบ Multi Controllers}
    \label{img:topo_cm}
\end{figure}

\pagebreak

โดยรูปภาพที่ \ref{img:topo_cs} คือ Topology แบบมี 1 \gls{controller} โดยจะทำการกำหนดให้\gls{controller}เป็นแบบ remote 
และรูปภาพที่ \ref{img:topo_cm} คือ Topology แบบมีหลาย\gls{controller} โดยจะทำการกำหนดให้\gls{controller}เป็นแบบ remote ทั้งหมด และกำหนด port เป็น 6633 6634 และ 6635
นอกจากนั้น เมื่อสร้าง script python ขึ้นมาแล้ว หรือดาวน์โหลดไฟล์จากลิงก์ที่ให้ไว้ จะทำการตั้งชื่อไฟล์เป็นดังนี้ Topology ที่มี 1 \gls{controller}
ดังรูปภาพที่ \ref{img:topo_cs} จะทำการตั้งชื่อเป็น single-topo.py และ Topology ที่มีหลาย\gls{controller} ดังรูปภาพที่ \ref{img:topo_cm} จะทำการตั้งชื่อเป็น multi-topo.py

ในการทดลองครั้งนี้ผู้ทำการทดลองได้ออกแบบการทดลองให้มีความเข้มงวดและรัดกุมมากที่สุด เพื่อให้ผลลัพธ์ที่ได้มีความคลาดเคลื่อนน้อยที่สุด
โดยผู้ทำการทดลองจะออกแบบการทดลองโดยให้มี 2 รูปแบบ 4 เงื่อนไข ดังนี้


เงื่อนไขที่ 1. (A1) group table แบบ 1 \gls{controller}\\
โดยจะทำการรัน ryu-manager ด้วย script ไฟล์ที่ชื่อ ex7\_group\_tables.py 
และทำการรันไฟล์ชื่อ single-topo.py ด้วย python เพื่อจำลอง Topology แบบ 1 \gls{controller}


เงื่อนไขที่ 2. (A2) group table แบบหลาย\gls{controller} \\
โดยจะทำการรัน ryu-manager ด้วย script ไฟล์ที่ชื่อ ex7\_group\_tables.py 
และทำการรันไฟล์ชื่อ multi-topo.py ด้วย python เพื่อจำลอง Topology แบบหลาย\gls{controller}


เงื่อนไขที่ 3. (A3) proxy arp แบบ 1 \gls{controller}\\
โดยจะทำการรัน ryu-manager ด้วย script ไฟล์ที่ชื่อ ex8\_arp\_proxy.py
และทำการรันไฟล์ชื่อ single-topo.py ด้วย python เพื่อจำลอง Topology แบบ 1 \gls{controller}

เงื่อนไขที่ 4. (A4) proxy arp แบบหลาย\gls{controller} \\
โดยจะทำการรัน ryu-manager ด้วย script ไฟล์ที่ชื่อ ex8\_arp\_proxy.py 
และทำการรันไฟล์ชื่อ multi-topo.py ด้วย python เพื่อจำลอง Topology แบบหลาย\gls{controller}

รูปแบบที่ 1. คือ ทำการทดสอบโดยจับแพ็กเก็ตทีละจำนวน 100,000 250,000 500,000 750,000 และ 1,000,000 แพ็กเก็ต
จากนั้นนำมาวิเคราะห์เปรียบเทียบหาประสิทธิภาพด้านเวลาและการตรวจจับการโจมตี

รูปแบบที่ 2. คือ ทำการทดสอบโดยจำแพ็กเก็ตจำนวน 1,000,000 และ 2,000,000 แพ็กเก็ต
จากนั้นนำมาวิเคราะห์เปรียบเทียบหาประสิทธิภาพด้านเวลาและการตรวจจับการโจมตี
\\

สาเหตุที่มีการทดลองสองรูปแบบนั้นก็เพื่อให้สามารถตรวจสอบได้ว่าผลลัพธ์ที่ได้จากการทดสอบมีความถูกต้องน่าเชื่อถือ
เป็นต้นว่าหากผลการทดสองจากทั้งสองรูปแบบสอดคล้องกัน ก็สามารถเชื่อได้ว่าผลการทดลองนั้นถูกต้อง นอกจากนั้นผู้ทำการทดลองยังได้กำหนดให้ลำดับวิธีของทั้ง 2 รูปแบบการทดลองมีความต่างกันในบางขั้นตอน เพื่อเพิ่มความเชื่อมั่นและความถูกต้องมากขั้นไปอีก
\\
\subsection{ขั้นตอนลำดับวิธีทดลอง}
ในการทดลอง ขั้นตอนลำดับจะต่างกันในส่วนของรูปแบบการทดลองเท่านั้น โดยที่เงื่อนไขต่างๆ จะยังคงมีลำดับการทดลองที่เหมือนกัน กล่าวคือในเงื่อนไขการทดลองนั้นจะต่างกันก็เพียงแต่ไฟล์ที่ทำการรันเท่านั้น โดยที่ลำดับขั้นตอนจะยังคงเหมือนเดิม 
ในแต่ละเงิ่อนไขการทดลองจะทำการรันทดสอบ 3 รอบ และเลือกค่ากลางมาทำการวิเคราะห์ กล่าวคือ เลือกค่าที่ไม่น้อยที่สุดและไม่มากที่สุดจากค่าที่ได้จากการทดสอบ

\subsubsection*{รูปแบบการทดลองที่ 1 บน 1 \gls{controller}}

ในขั้นต้นผู้ทำการทดลองต้องทำการเปิด gnome terminal ขึ้นมาและเปิดหน้าต่างย่อยเป็นจำนวน 3 หน้าต่าง
ทุกหน้าต่างจะต้องทำการเชื่อมต่อไปยัง\gls{server}ด้วย ssh ให้เรียบร้อย
โดย หน้าต่างที่ 1 จะทำการรัน ryu-manager หน้าต่างที่ 2 จะทำการรัน mininet และหน้าต่างที่ 3 จะทำการรัน wireshark
โดยมีลำดับขั้นตอนการทดลองดังนี้


1. หน้าต่างที่ 1 ทำการรันคำสั่ง    
\$ ryu-manager learn-sdn-with-ryu/ryu-exercises/<python script> \\
โดย <python script> จะแทนด้วย ex7\_group\_tables.py หรือ ex8\_arp\_proxy.py ตามเงื่อนไขการทดลอง

2. หน่าต่างที่ 2 ทำการรันคำสั่ง    
\$ sudo python3 single\_topo.py \\
เพื่อทำการจำลอง topology แบบ 1 \gls{controller} ขึ้นมา

3. หน้าต่างที่ 3 ทำการรันคำสั่ง 
\$ sudo wireshark \\
โดยทำการจับแพ็กเก็ตที่ s4-eth1 (ที่เชื่อมกับ h2) และ s11-eth2 (ที่เชื่อมกับ h12) เป็นจำนวนที่เคยกล่าวไปข้างต้น

4. หลังจากที่ wireshark ทำการเริ่มจับแพ็กเก็ต หน้าต่างที่ 2 จะทำการรันคำสั่ง \$ mininet> h2 hping3 h12 -S --flood -V
โดยทันที เพื่อทำการให้ h2 โจมตีไปที่ h12

5. หลังจากทำการจับแพ็กเก็ตครบตามจำนวนที่กำหนด หน้าต่างที่ 2 ทำการหยุดการโจมตี จากนั้นทำการเก็บค่าสถิติ และทำการปิดโปรแกรม wireshark
หน้าต่างที่ 2 ทำการออกจาก mininet และรันคำสั่ง \$ sudo mn -c หากมีการทดสอบรอบต่อไปต้องเริ่มจากขั้นตอนที่ 1 ใหม่เท่านั้น

\subsubsection*{รูปแบบการทดลองที่ 1 บนหลาย\gls{controller}}

ในขั้นต้นผู้ทำการทดลองต้องทำการเปิด gnome terminal ขึ้นมาและเปิดหน้าต่างย่อยเป็นจำนวน 5 หน้าต่าง
ทุกหน้าต่างจะต้องทำการเชื่อมต่อไปยัง\gls{server}ด้วย ssh ให้เรียบร้อย
โดย หน้าต่างที่ 1 2 และ 3 จะทำการรัน ryu-manager หน้าต่างที่ 4 จะทำการรัน mininet และหน้าต่างที่ 5 จะทำการรัน wireshark
โดยมีลำดับขั้นตอนการทดลองดังนี้

1. หน้าต่างที่ 1 ทำการรันคำสั่ง \\
\$ ryu-manager learn-sdn-with-ryu/ryu-exercises/<python script> --ofp-tcp-listen-port 6633 \\
\indent หน้าต่างที่ 2 ทำการรันคำสั่ง \\
\$ ryu-manager learn-sdn-with-ryu/ryu-exercises/<python script> --ofp-tcp-listen-port 6634  \\
\indent หน้าต่างที่ 3 ทำการรันคำสั่ง \\
\$ ryu-manager learn-sdn-with-ryu/ryu-exercises/<python script> --ofp-tcp-listen-port 6635  \\
โดย <python script> จะแทนด้วย 
ex7\_group\_tables.py หรือ ex8\_arp\_proxy.py
ตามเงื่อนไขการทดลอง

2. หน่าต่างที่ 4 ทำการรันคำสั่ง    
\$ sudo python3 multi\_topo.py \\
เพื่อทำการจำลอง topology แบบหลาย\gls{controller} ขึ้นมา

3. หน้าต่างที่ 5 ทำการรันคำสั่ง 
\$ sudo wireshark 
โดยทำการจับแพ็กเก็ตที่ s4-eth1 (ที่เชื่อมกับ h2) และ s11-eth2 (ที่เชื่อมกับ h12) เป็นจำนวนที่เคยกล่าวไปข้างต้น

4. หลังจากที่ wireshark ทำการเริ่มจับแพ็กเก็ต หน้าต่างที่ 4 จะทำการรันคำสั่ง \$ mininet> h2 hping3 h12 -S --flood -V
โดยทันที เพื่อทำการให้ h2 โจมตีไปที่ h12

5. หลังจากทำการจับแพ็กเก็ตครบตามจำนวนที่กำหนด หน้าต่างที่ 4 ทำการหยุดการโจมตี จากนั้นทำการเก็บค่าสถิติ และทำการปิดโปรแกรม wireshark
หน้าต่างที่ 2 ทำการออกจาก mininet และรันคำสั่ง \$ sudo mn -c หากมีการทดสอบรอบต่อไปต้องเริ่มจากขั้นตอนที่ 1 ใหม่เท่านั้น

\subsubsection*{รูปแบบการทดลองที่ 2 บน 1 \gls{controller}}

ในขั้นต้นผู้ทำการทดลองต้องทำการเปิด gnome terminal ขึ้นมาและเปิดหน้าต่างย่อยเป็นจำนวน 3 หน้าต่าง
ทุกหน้าต่างจะต้องทำการเชื่อมต่อไปยัง\gls{server}ด้วย ssh ให้เรียบร้อย
โดย หน้าต่างที่ 1 จะทำการรัน ryu-manager หน้าต่างที่ 2 จะทำการรัน mininet และหน้าต่างที่ 3 จะทำการรัน wireshark
โดยมีลำดับขั้นตอนการทดลองดังนี้


1. หน้าต่างที่ 1 ทำการรันคำสั่ง    
\$ ryu-manager learn-sdn-with-ryu/ryu-exercises/<python script> \\
โดย <python script> จะแทนด้วย ex7\_group\_tables.py หรือ ex8\_arp\_proxy.py ตามเงื่อนไขการทดลอง

2. หน่าต่างที่ 2 ทำการรันคำสั่ง    
\$ sudo python3 single\_topo.py \\
เพื่อทำการจำลอง topology แบบ 1 \gls{controller} ขึ้นมา

3. หน้าต่างที่ 2 จะทำการรันคำสั่ง \$ mininet> h2 hping3 h12 -S --flood -V
เพื่อทำการให้ h2 โจมตีไปที่ h12

4. หน้าต่างที่ 3 ทำการรันคำสั่ง 
\$ sudo wireshark 
จากนั้นจะทำการเริ่มจับแพ็กเก็ตหลังจากเริ่มการโจมตีแล้ว 30 วินาที
โดยทำการจับแพ็กเก็ตที่ s4-eth1 (ที่เชื่อมกับ h2) และ s11-eth2 (ที่เชื่อมกับ h12) เป็นจำนวนที่เคยกล่าวไปข้างต้น


5. หลังจากทำการจับแพ็กเก็ตครบตามจำนวนที่กำหนด หน้าต่างที่ 2 ทำการหยุดการโจมตี จากนั้นทำการเก็บค่าสถิติ และทำการปิดโปรแกรม wireshark
หน้าต่างที่ 2 ทำการออกจาก mininet และรันคำสั่ง \$ sudo mn -c หากมีการทดสอบรอบต่อไปต้องเริ่มจากขั้นตอนที่ 1 ใหม่เท่านั้น

\subsubsection*{รูปแบบการทดลองที่ 2 บนหลาย\gls{controller}}

ในขั้นต้นผู้ทำการทดลองต้องทำการเปิด gnome terminal ขึ้นมาและเปิดหน้าต่างย่อยเป็นจำนวน 5 หน้าต่าง
ทุกหน้าต่างจะต้องทำการเชื่อมต่อไปยัง\gls{server}ด้วย ssh ให้เรียบร้อย
โดย หน้าต่างที่ 1 2 และ 3 จะทำการรัน ryu-manager หน้าต่างที่ 4 จะทำการรัน mininet และหน้าต่างที่ 5 จะทำการรัน wireshark
โดยมีลำดับขั้นตอนการทดลองดังนี้

1. หน้าต่างที่ 1 ทำการรันคำสั่ง \\
\$ ryu-manager learn-sdn-with-ryu/ryu-exercises/<python script> --ofp-tcp-listen-port 6633 \\
\indent หน้าต่างที่ 2 ทำการรันคำสั่ง \\
\$ ryu-manager learn-sdn-with-ryu/ryu-exercises/<python script> --ofp-tcp-listen-port 6634  \\
\indent หน้าต่างที่ 3 ทำการรันคำสั่ง \\
\$ ryu-manager learn-sdn-with-ryu/ryu-exercises/<python script> --ofp-tcp-listen-port 6635  \\
โดย <python script> จะแทนด้วย 
ex7\_group\_tables.py หรือ ex8\_arp\_proxy.py
ตามเงื่อนไขการทดลอง

2. หน่าต่างที่ 4 ทำการรันคำสั่ง    
\$ sudo python3 multi\_topo.py เพื่อทำการจำลอง topology แบบหลาย\gls{controller} ขึ้นมา

4. หน้าต่างที่ 4 จะทำการรันคำสั่ง \$ mininet> h2 hping3 h12 -S --flood -V
เพื่อทำการให้ h2 โจมตีไปที่ h12

3. หน้าต่างที่ 5 ทำการรันคำสั่ง 
\$ sudo wireshark 
จากนั้นจะทำการเริ่มจับแพ็กเก็ตหลังจากเริ่มการโจมตีแล้ว 30 วินาที
โดยทำการจับแพ็กเก็ตที่ s4-eth1 (ที่เชื่อมกับ h2) และ s11-eth2 (ที่เชื่อมกับ h12) เป็นจำนวนที่เคยกล่าวไปข้างต้น

5. หลังจากทำการจับแพ็กเก็ตครบตามจำนวนที่กำหนด หน้าต่างที่ 4 ทำการหยุดการโจมตี จากนั้นทำการเก็บค่าสถิติ และทำการปิดโปรแกรม wireshark
หน้าต่างที่ 2 ทำการออกจาก mininet และรันคำสั่ง \$ sudo mn -c หากมีการทดสอบรอบต่อไปต้องเริ่มจากขั้นตอนที่ 1 ใหม่เท่านั้น
\\\\

\section{ผลลัพธ์และการอภิปราย (Results and discussion)}
เหตุผลที่ต้องทำการรัน 3 รอบ ในแต่ละเงื่อนไข เนื่องมาจากในขั้นตอนการเตรียมเครื่องมือและซอฟต์แวร์นั้นได้มีการรันทดสอบ และพบว่าค่ามีความแกว่งในบางครั้ง โดยในการรันครั้งแรกดรอปแพ็กเกจได้ 8\% ในการรันครั้งที่ 2 อาจจะเพิ่มเป็น 30\% หรือ 35\% 
ดังนั้นจึงได้มีการออกแบบขั้นตอนการทดลองที่เหมือนกันมากที่สุดเพื่อลดปัจจัยที่อาจจะกระทบและทำให้ผลลัพธ์เปลี่ยนไปอย่างมีนัยสำคัญ นอกจากนั้นจึงได้มีการรันทดสอบ 3 รอบของแต่ละรูปแบบเพื่อลดโอกาสที่ค่าจะออกมาคลาดเคลื่อนหรือผิดไปจากที่ควรจะเป็น 
และในการรันทดสอบนั้นมีการจับแพ็กเก็ตด้วย wireshark ที่ s4-eth1 และ s11-eth2 ซึ่งทำให้มี 2 ค่าที่เกิดขึ้นจากทั้ง 2 ขาของการจับแพ็กเกจ ทางผู้ทดลองต้องการทำให้เป็นค่าเพียงหนึ่งค่า
จึงจะทำการรวมค่าจาก s4-eth1 และ s11-eth2 ให้เป็นค่าหนึ่งโดยการบวก และในการรันทดสอบ 3 ครั้งนั้น 
ผู้ทำการทดลองพิจารณาแล้วว่าจะเลือกค่าอัตราการดรอปแพ็กเกจที่เป็นค่ากลางมาใช้ในการวิเคราะห์
\\
\subsection{ผลลัพธ์การทดลอง}
\subsubsection*{ผลลัพธ์ด้านเวลา}

%%%%%
%
%% Here
%
%%%%%
ในการทดสอบรูปแบบที่ 1 มีการจับแพ็กเก็ตหลายระดับ จึงอาจจะทำให้แพ็กเก็จจำนวนมากกว่าจะใช้เวลามากกว่า
จึงได้ใช้สูตรคำนวนเพื่อให้สามารถเทียบวัดได้ โดยใช้สูตร
$$P=t \times \frac{m}{d}$$

โดยที่ P คือค่าคะแนน t คือเวลาที่ได้จากการวัด m คือจำนวนแพ็กเก็ตมากที่สุดที่ทำการวัด d คือจำนวนแพ็กเก็ตที่วัดในรอบนั้น

\begin{figure}[h!]
    \centering
    \begin{tikzpicture}
        \begin{axis}[
            title=\textbf{ประสิทธิภาพด้านเวลา การทดลองรูปแบบที่ 1},
             width=0.8\textwidth,
            height=0.5\textwidth,
            x tick label style={
		    /pgf/number format/1000 sep=},
            legend style={at={(1,1)},
	        anchor=north east},
            xlabel=packet,
            ylabel=P (ยิ่งน้อยยิ่งดี),
            % ybar interval=0.7,
            % ybar,
            ymax=50,
            ymin=0,
            symbolic x coords={
            100000,
            250000,
            500000,
            750000,
            1000000
            },
            xtick=data,
            enlarge x limits=0.15,
            ]
            \addplot coordinates {
                % group table single
                (100000,26.89)
                (250000,20.956)
                (500000,21.614)
                (750000,18.601)
                (1000000,20.595)
            };
            \addplot coordinates {
                % group table multiple
                (100000,19.75)
                (250000,8.724)
                (500000,14.478)
                (750000,33.437)
                (1000000,33.052)
            };
            \addplot coordinates {
                % arp proxy single
                (100000,28.49)
                (250000,20.104)
                (500000,21.776)
                (750000,20.478)
                (1000000,21.294)
            };
             \addplot coordinates {
                % arp proxy multiple
                (100000,42.3)
                (250000,35.8)
                (500000,26.944)
                (750000,29.342)
                (1000000,21.489)
            };
           \legend{
            A1,A2,A3,A4
          }
        \end{axis}
    \end{tikzpicture}
    \caption{แผนภูมิแสดงผลลัพธ์ประสิทธิภาพด้านเวลาของการทดลองรูปแบบที่ 1}
    \label{img:graph_time_1}
\end{figure}

จากรูปภาพที่ \ref{img:graph_time_1} (A1) มีประสิทธิภาพไม่ต่างกันอย่างมีนัยสำคัญเมื่อแพ็กเก็ตเพิ่มขึ้น แต่มีแนวโน้มดีขึ้นเมื่อแพ็กเก็ตเพิ่มขึ้น
(A2) มีแนวโน้มมีประสิทธิภาพน้อยลงเมื่อแพ็กเก็ตเพิ่มขึ้น แต่ไม่สม่ำเสมอ 
(A3) ประสิทธิภาพดีขึ้นเล็กน้อยอย่างไม่มีนัยสำคัญเมื่อแพ็กเก็ตเพิ่มขึ้น 
(A4) มีประสิทธิภาพดีขึ้นเมื่อแพ็กเก็ตเพิ่มขึ้น และมีแนวโน้มประสิทธิภาพจะดีขึ้นอีก เมื่อแพ็กเก็ตเพิ่มขึ้น

\begin{figure}[h!]
    \centering
    \begin{tikzpicture}
        \begin{axis}[
            title=\textbf{ประสิทธิภาพด้านเวลา การทดลองรูปแบบที่ 2},
             width=0.8\textwidth,
            height=0.5\textwidth,
            x tick label style={
		    /pgf/number format/1000 sep=},
            legend style={at={(1,1)},
	        anchor=north east},
            xlabel=packet,
            ylabel=P (ยิ่งน้อยยิ่งดี),
            % ybar interval=0.7,
            % ybar,
            ymax=50,
            ymin=0,
            symbolic x coords={
            1000000,
            2000000
            },
            xtick=data,
            enlarge x limits=0.15,
            ]
            \addplot coordinates {
                % group table single
                (1000000,27.994)
                (2000000,21.84)
            };
            \addplot coordinates {
                % group table multiple
                (1000000,29.55)
                (2000000,13.319)
            };
            \addplot coordinates {
                % arp proxy single
                (1000000,21.66)
                (2000000,23.125)
            };
             \addplot coordinates {
                % arp proxy multiple
                (1000000,18.478)
                (2000000,17.851)
            };
           \legend{
            A1,A2,A3,A4
          }
        \end{axis}
    \end{tikzpicture}
    \caption{แผนภูมิแสดงผลลัพธ์ประสิทธิภาพด้านเวลาของการทดลองรูปแบบที่ 2}
    \label{img:graph_time_2}
\end{figure}

จากรูปภาพที่ \ref{img:graph_time_2} (A1) มีประสิทธิภาพดีขึ้นเมื่อแพ็กเก็ตเพิ่มขึ้น
(A2) มีประสิทธิภาพดีขึ้นอย่างเห็นได้ชัดเมื่อแพ็กเก็ตเพิ่มขึ้น
(A3) ประสิทธิภาพเพิ่มขึ้นเล็กน้อย แต่ไม่มีนัยสำคัญ
(A4) ไม่มีการเปลี่ยนแปลงอย่างเห็นได้ชัด อาจกล่าวได้ว่าประสิทธิภาพเท่าเดิม

\subsubsection*{ผลลัพธ์ประสิทธิภาพด้านการตรวจจับการโจมตี}

% attack table 
\begin{figure}[h!]
    \centering
    \begin{tikzpicture}
        \begin{axis}[
            title=\textbf{ประสิทธิการตรวจจับการโจมตี การทดลองรูปแบบที่ 1},
             width=0.8\textwidth,
            height=0.5\textwidth,
            x tick label style={
		    /pgf/number format/1000 sep=},
            legend style={at={(1,1)},
	        anchor=north east},
            xlabel=packet,
            ylabel=droped packet (\%),
            % ybar interval=0.7,
            % ybar,
            ymax=30,
            ymin=0,
            symbolic x coords={
            100000,
            250000,
            500000,
            750000,
            1000000
            },
            xtick=data,
            enlarge x limits=0.15,
            ]
            \addplot coordinates {
                % group table single
                (100000,21.6)
                (250000,17.2)
                (500000,14.9)
                (750000,5.8)
                (1000000, 11.3)
            };
            \addplot coordinates {
                % group table multiple
                (100000,9.8)
                (250000,4.4)
                (500000,8)
                (750000,4.5)
                (1000000,4)
            };
            \addplot coordinates {
                % arp proxy single
                (100000,20.9)
                (250000,10.7)
                (500000,7.5)
                (750000,7.7)
                (1000000,8)
            };
             \addplot coordinates {
                % arp proxy multiple
                (100000,8.1)
                (250000,10.9)
                (500000,5.4)
                (750000,9.3)
                (1000000,6.1)
            };
           \legend{
            A1,A2,A3,A4
          }
        \end{axis}
    \end{tikzpicture}
    \caption{แผนภูมิแสดงผลลัพธ์ประสิทธิภาพด้านการตรวจจับการโจมตีของการทดลองรูปแบบที่ 1}
    \label{img:graph_atk_1}
\end{figure}

จากรูปที่ \ref*{img:graph_atk_1} จะเห็นว่า (A1) มีประสิทธิภาพการตรวจจับที่น้อยลงเมื่อแพ็กเก็ตเพิ่มขึ้น
(A2) มีประสิทธิภาพไม่คงที่ในแต่ละจำนวนแพ็กเก็ต และไม่มีความต่างอย่างมีนัยสำคัญ
(A3) มีประสิทธิภาพน้อยลงเมื่อแพ็กเก็ตเพิ่มขึ้น แต่เมื่อถึงจุดหนึ่งประสิทธิภาพจะไม่ลดลงอีก
(A4) มีประสิทธิภาพไม่คงที่ และไม่มีแนวโน้มจะลดลงหรือเพิ่มขึ้น

\begin{figure}[h!]
    \centering
    \begin{tikzpicture}
        \begin{axis}[
            title=\textbf{ประสิทธิการตรวจจับการโจมตี การทดลองรูปแบบที่ 2},
             width=0.8\textwidth,
            height=0.5\textwidth,
            x tick label style={
		    /pgf/number format/1000 sep=},
            legend style={at={(1,1)},
	        anchor=north east},
            xlabel=packet,
            ylabel=droped packet (\%),
            % ybar interval=0.7,
            % ybar,
            ymax=30,
            ymin=0,
            symbolic x coords={
            1000000,
            2000000
            },
            xtick=data,
            enlarge x limits=0.15,
            ]
            \addplot coordinates {
                % group table single
                (1000000, 8.1)
                (2000000, 10.5)
            };
            \addplot coordinates {
                % group table multiple
                (1000000,12.5)
                (2000000,21)
            };
            \addplot coordinates {
                % arp proxy single
                (1000000,12.4)
                (2000000,16.5)
            };
             \addplot coordinates {
                % arp proxy multiple
                (1000000,20.3)
                (2000000, 12.4)
            };
           \legend{
            A1,A2,A3,A4
          }
        \end{axis}
    \end{tikzpicture}
    \caption{แผนภูมิแสดงผลลัพธ์ประสิทธิภาพด้านการตรวจจับการโจมตีของการทดลองรูปแบบที่ 2}
    \label{img:graph_atk_2}
\end{figure}

จากรูปภาพที่ \ref{img:graph_atk_2} (A1) มีประสิทธิภาพเพิ่มขึ้นเมื่อแพ็กเก็ตเพิ่มขึ้น และมีแนวโน้มประสิทธิภาพเพิ่มขึ้น
(A2) ประสิทธิภาพเพิ่มขึ้นอย่างเห็นได้ชัด และมีแนวโน้มอย่างสูงว่าประสิทธิภาพจะเพิ่มขึ้นอีกหากแพ็กเก็ตเพิ่มขึ้น
(A3) มีประสิทธิภาพเพิ่มขึ้นเมื่อแพ็กเก็ตเพิ่มขึ้น มีแนวโน้มเล็กน้อยว่าประสิทธิภาพจะเพิ่มขึ้นอีก
(A4) มีประสิทธิภาพลดลงเมื่อแพ็กเก็ตเพิ่มขึ้น ซึ่งสวนทางกับเงื่อนไขอื่นอย่างเห็นได้ชัด
\\
% \clearpage
\subsection{สรุปผลการทดลอง}

จากที่ได้กล่าวไปว่า การที่ต้องทดลอง 2 รูปแบบ ก็เพื่อที่จะเปรียบเทียบและเพิ่มความถูกต้อง
เป็นต้นว่า หากทั้งสองรูปแบบสอดคล้องกัน ก็เป็นที่น่าเชื่อว่าผลการทดลองนั้นออกมาถูกต้อง
แต่หากทั้งสองรูปแบบออกมาขัดแย้งกัน ก็เป็นไปได้ว่าการทดลองนี้ไม่ถูกต้อง หรือมีข้อผิดพลาด
ดังนั้น ในการที่จะสามารถบอกได้ชัดเจนเกี่ยวกับผลทดลองนั้น จำเป็นต้องดูทั้ง 2 รูปแบบประกอบกัน

โดยในประสิทธิภาพด้านเวลานั้น จะเห็นได้อย่างชัดเจนว่า (A2) นั้นผลออกมาขัดแย้งกันอย่างชัดเจน จึงกล่าวได้ว่า ผลลัพธ์การทดสอบของ (A2) นั้นมีความน่าเชื่อถือต่ำ
หรืออาจจะมีข้อผิดพลาดในการทดลอง จึงไม่อาจจะนำ (A2) ไปใช้ในการวิเคราะห์อื่นๆ ได้
ส่วนในผลลัพธ์อื่นๆ นั้น จะเห็นว่าทั้ง 2 รูปแบบไม่ได้ขัดแย้งกันอย่างชัดเจน ผลลัพธ์จึงมีความน่าเชื่อถือ
โดย (A4) นั้นมีประสิทธิภาพที่ดีขึ้นเมื่อแพ็กเก็ตมีจำนวนเพิ่มขึ้น ส่วนที่เหลือนั้นไม่ได้มีความต่างอย่างชัดเจน
สามารถกล่าวได้ว่า (A4) ให้ประสิทธิภาพด้านเวลาดีที่สุด

ในประสิทธิภาพด้านการตรวจจับการทดลองนั้น จะเห็นได้อย่างชัดเจนเช่นเดียวกับประสิทธิภาพด้านเวลา ว่า (A2) นั้นมีผลลัพธ์ที่ขัดแย้งกันจากทัง 2 รูปแบบ ดังนั้น (A2) จึงมีความน่าเชื่อถือต่ำ
ในส่วน (A1) นั้น แม้ดูมีแนวโน้มว่าจะขัดแย้งกัน แต่เมื่อพิจารณาผลลัพธ์ในรูปแบบที่ 2 แล้ว ค่ามีการเพิ่มขึ้นเล็กน้อยเพียงเท่านั้น ทำให้ยังไม่ได้มีความขัดแย้งอย่างเห็นได้ชัด
ส่วนอื่นๆ นั้น ถือว่าค่อนข้างมีความสอดคล้องกัน โดยที่ไม่ได้มีข้อใดมีประสิทธิภาพที่ดีกว่าอย่างเห็นได้ชัด หากดูตามกราฟแล้ว พอจะพูดได้ว่า (A3) ให้ประสิทธิภาพโดยรวมที่ดีที่สุด

ประสิทธิภาพด้านการตรวจจับการโจมตีนั้น แม้จะไม่ได้มีความขัดแย้งกันอย่างชัดเจน แต่ผลลัพธ์ก็ไม่ได้เป็นไปในทิศทางเดียวกันมากนัก ดังนั้นจึงต้องมีความระมัดระวังมากเป็นพิเศษในการนำผลลัพธ์ด้านการตรวจจับแพ็กเก็ตไปวิเคราะห์ โดยหากต้องการนำข้อมูลใดๆ ไปวิเคราะห์หรือประกอบการตัดสินใจ แนะนำให้พิจารณาเลือกรูปแบบที่ 2 เป็นอันดับแรก
เพราะมีช่วงจำนวนการจับแพ็กเก็ตมากกว่า
\\
\subsection{ข้อเสนอแนะ}

ในการทดลองนี้ แม้ผู้ทำการทดลองจะออกแบบการทดลองให้มีความรัดกุม และเข้มงวดมากที่สุด
แต่ก็หลีกเลี่ยงไม่ได้ที่จะมีข้อผิดพลาด หรือผลลัพธ์ที่คลาดเคลื่อน แม้แต่ผลลัพธ์ที่ไม่สอดคล้องเองก็ตาม
อย่างไรก็ดี หากต้องการนำผลการทดลองไปใช้ ควรทำการทดสอบและวัดผลด้วยตัวเอง ตามวิธีที่เขียนไว้ หรืออาจจะวิธีการอื่นที่ดีกว่า ซึ่งอาจจะเพิ่มความรัดกุมขึ้นอีก
หรือเพิ่มความแม่นยำของผลลัพธ์ โดยอิงตามทฤษฎีมากขึ้น และควบคุมตัวแปรได้อย่างเรียบร้อย
% \section{ผลลัพธ์และการอภิปราย (Results and discussion)}
เหตุผลที่ต้องทำการรัน 3 รอบ ในแต่ละเงื่อนไข เนื่องมาจากในขั้นตอนการเตรียมเครื่องมือและซอฟต์แวร์นั้นได้มีการรันทดสอบ และพบว่าค่ามีความแกว่งในบางครั้ง โดยในการรันครั้งแรกดรอปแพ็กเกจได้ 8\% ในการรันครั้งที่ 2 อาจจะเพิ่มเป็น 30\% หรือ 35\% 
ดังนั้นจึงได้มีการออกแบบขั้นตอนการทดลองที่เหมือนกันมากที่สุดเพื่อลดปัจจัยที่อาจจะกระทบและทำให้ผลลัพธ์เปลี่ยนไปอย่างมีนัยสำคัญ นอกจากนั้นจึงได้มีการรันทดสอบ 3 รอบของแต่ละรูปแบบเพื่อลดโอกาสที่ค่าจะออกมาคลาดเคลื่อนหรือผิดไปจากที่ควรจะเป็น 
และในการรันทดสอบนั้นมีการจับแพ็กเก็ตด้วย wireshark ที่ s4-eth1 และ s11-eth2 ซึ่งทำให้มี 2 ค่าที่เกิดขึ้นจากทั้ง 2 ขาของการจับแพ็กเกจ ทางผู้ทดลองต้องการทำให้เป็นค่าเพียงหนึ่งค่า
จึงจะทำการรวมค่าจาก s4-eth1 และ s11-eth2 ให้เป็นค่าหนึ่งโดยการบวก และในการรันทดสอบ 3 ครั้งนั้น 
ผู้ทำการทดลองพิจารณาแล้วว่าจะเลือกค่าอัตราการดรอปแพ็กเกจที่เป็นค่ากลางมาใช้ในการวิเคราะห์
\\
\subsection{ผลลัพธ์การทดลอง}
\subsubsection*{ผลลัพธ์ด้านเวลา}

%%%%%
%
%% Here
%
%%%%%
ในการทดสอบรูปแบบที่ 1 มีการจับแพ็กเก็ตหลายระดับ จึงอาจจะทำให้แพ็กเก็จจำนวนมากกว่าจะใช้เวลามากกว่า
จึงได้ใช้สูตรคำนวนเพื่อให้สามารถเทียบวัดได้ โดยใช้สูตร
$$P=t \times \frac{m}{d}$$

โดยที่ P คือค่าคะแนน t คือเวลาที่ได้จากการวัด m คือจำนวนแพ็กเก็ตมากที่สุดที่ทำการวัด d คือจำนวนแพ็กเก็ตที่วัดในรอบนั้น

\begin{figure}[h!]
    \centering
    \begin{tikzpicture}
        \begin{axis}[
            title=\textbf{ประสิทธิภาพด้านเวลา การทดลองรูปแบบที่ 1},
             width=0.8\textwidth,
            height=0.5\textwidth,
            x tick label style={
		    /pgf/number format/1000 sep=},
            legend style={at={(1,1)},
	        anchor=north east},
            xlabel=packet,
            ylabel=P (ยิ่งน้อยยิ่งดี),
            % ybar interval=0.7,
            % ybar,
            ymax=50,
            ymin=0,
            symbolic x coords={
            100000,
            250000,
            500000,
            750000,
            1000000
            },
            xtick=data,
            enlarge x limits=0.15,
            ]
            \addplot coordinates {
                % group table single
                (100000,26.89)
                (250000,20.956)
                (500000,21.614)
                (750000,18.601)
                (1000000,20.595)
            };
            \addplot coordinates {
                % group table multiple
                (100000,19.75)
                (250000,8.724)
                (500000,14.478)
                (750000,33.437)
                (1000000,33.052)
            };
            \addplot coordinates {
                % arp proxy single
                (100000,28.49)
                (250000,20.104)
                (500000,21.776)
                (750000,20.478)
                (1000000,21.294)
            };
             \addplot coordinates {
                % arp proxy multiple
                (100000,42.3)
                (250000,35.8)
                (500000,26.944)
                (750000,29.342)
                (1000000,21.489)
            };
           \legend{
            A1,A2,A3,A4
          }
        \end{axis}
    \end{tikzpicture}
    \caption{แผนภูมิแสดงผลลัพธ์ประสิทธิภาพด้านเวลาของการทดลองรูปแบบที่ 1}
    \label{img:graph_time_1}
\end{figure}

จากรูปภาพที่ \ref{img:graph_time_1} (A1) มีประสิทธิภาพไม่ต่างกันอย่างมีนัยสำคัญเมื่อแพ็กเก็ตเพิ่มขึ้น แต่มีแนวโน้มดีขึ้นเมื่อแพ็กเก็ตเพิ่มขึ้น
(A2) มีแนวโน้มมีประสิทธิภาพน้อยลงเมื่อแพ็กเก็ตเพิ่มขึ้น แต่ไม่สม่ำเสมอ 
(A3) ประสิทธิภาพดีขึ้นเล็กน้อยอย่างไม่มีนัยสำคัญเมื่อแพ็กเก็ตเพิ่มขึ้น 
(A4) มีประสิทธิภาพดีขึ้นเมื่อแพ็กเก็ตเพิ่มขึ้น และมีแนวโน้มประสิทธิภาพจะดีขึ้นอีก เมื่อแพ็กเก็ตเพิ่มขึ้น

\begin{figure}[h!]
    \centering
    \begin{tikzpicture}
        \begin{axis}[
            title=\textbf{ประสิทธิภาพด้านเวลา การทดลองรูปแบบที่ 2},
             width=0.8\textwidth,
            height=0.5\textwidth,
            x tick label style={
		    /pgf/number format/1000 sep=},
            legend style={at={(1,1)},
	        anchor=north east},
            xlabel=packet,
            ylabel=P (ยิ่งน้อยยิ่งดี),
            % ybar interval=0.7,
            % ybar,
            ymax=50,
            ymin=0,
            symbolic x coords={
            1000000,
            2000000
            },
            xtick=data,
            enlarge x limits=0.15,
            ]
            \addplot coordinates {
                % group table single
                (1000000,27.994)
                (2000000,21.84)
            };
            \addplot coordinates {
                % group table multiple
                (1000000,29.55)
                (2000000,13.319)
            };
            \addplot coordinates {
                % arp proxy single
                (1000000,21.66)
                (2000000,23.125)
            };
             \addplot coordinates {
                % arp proxy multiple
                (1000000,18.478)
                (2000000,17.851)
            };
           \legend{
            A1,A2,A3,A4
          }
        \end{axis}
    \end{tikzpicture}
    \caption{แผนภูมิแสดงผลลัพธ์ประสิทธิภาพด้านเวลาของการทดลองรูปแบบที่ 2}
    \label{img:graph_time_2}
\end{figure}

จากรูปภาพที่ \ref{img:graph_time_2} (A1) มีประสิทธิภาพดีขึ้นเมื่อแพ็กเก็ตเพิ่มขึ้น
(A2) มีประสิทธิภาพดีขึ้นอย่างเห็นได้ชัดเมื่อแพ็กเก็ตเพิ่มขึ้น
(A3) ประสิทธิภาพเพิ่มขึ้นเล็กน้อย แต่ไม่มีนัยสำคัญ
(A4) ไม่มีการเปลี่ยนแปลงอย่างเห็นได้ชัด อาจกล่าวได้ว่าประสิทธิภาพเท่าเดิม

\subsubsection*{ผลลัพธ์ประสิทธิภาพด้านการตรวจจับการโจมตี}

% attack table 
\begin{figure}[h!]
    \centering
    \begin{tikzpicture}
        \begin{axis}[
            title=\textbf{ประสิทธิการตรวจจับการโจมตี การทดลองรูปแบบที่ 1},
             width=0.8\textwidth,
            height=0.5\textwidth,
            x tick label style={
		    /pgf/number format/1000 sep=},
            legend style={at={(1,1)},
	        anchor=north east},
            xlabel=packet,
            ylabel=droped packet (\%),
            % ybar interval=0.7,
            % ybar,
            ymax=30,
            ymin=0,
            symbolic x coords={
            100000,
            250000,
            500000,
            750000,
            1000000
            },
            xtick=data,
            enlarge x limits=0.15,
            ]
            \addplot coordinates {
                % group table single
                (100000,21.6)
                (250000,17.2)
                (500000,14.9)
                (750000,5.8)
                (1000000, 11.3)
            };
            \addplot coordinates {
                % group table multiple
                (100000,9.8)
                (250000,4.4)
                (500000,8)
                (750000,4.5)
                (1000000,4)
            };
            \addplot coordinates {
                % arp proxy single
                (100000,20.9)
                (250000,10.7)
                (500000,7.5)
                (750000,7.7)
                (1000000,8)
            };
             \addplot coordinates {
                % arp proxy multiple
                (100000,8.1)
                (250000,10.9)
                (500000,5.4)
                (750000,9.3)
                (1000000,6.1)
            };
           \legend{
            A1,A2,A3,A4
          }
        \end{axis}
    \end{tikzpicture}
    \caption{แผนภูมิแสดงผลลัพธ์ประสิทธิภาพด้านการตรวจจับการโจมตีของการทดลองรูปแบบที่ 1}
    \label{img:graph_atk_1}
\end{figure}

จากรูปที่ \ref*{img:graph_atk_1} จะเห็นว่า (A1) มีประสิทธิภาพการตรวจจับที่น้อยลงเมื่อแพ็กเก็ตเพิ่มขึ้น
(A2) มีประสิทธิภาพไม่คงที่ในแต่ละจำนวนแพ็กเก็ต และไม่มีความต่างอย่างมีนัยสำคัญ
(A3) มีประสิทธิภาพน้อยลงเมื่อแพ็กเก็ตเพิ่มขึ้น แต่เมื่อถึงจุดหนึ่งประสิทธิภาพจะไม่ลดลงอีก
(A4) มีประสิทธิภาพไม่คงที่ และไม่มีแนวโน้มจะลดลงหรือเพิ่มขึ้น

\begin{figure}[h!]
    \centering
    \begin{tikzpicture}
        \begin{axis}[
            title=\textbf{ประสิทธิการตรวจจับการโจมตี การทดลองรูปแบบที่ 2},
             width=0.8\textwidth,
            height=0.5\textwidth,
            x tick label style={
		    /pgf/number format/1000 sep=},
            legend style={at={(1,1)},
	        anchor=north east},
            xlabel=packet,
            ylabel=droped packet (\%),
            % ybar interval=0.7,
            % ybar,
            ymax=30,
            ymin=0,
            symbolic x coords={
            1000000,
            2000000
            },
            xtick=data,
            enlarge x limits=0.15,
            ]
            \addplot coordinates {
                % group table single
                (1000000, 8.1)
                (2000000, 10.5)
            };
            \addplot coordinates {
                % group table multiple
                (1000000,12.5)
                (2000000,21)
            };
            \addplot coordinates {
                % arp proxy single
                (1000000,12.4)
                (2000000,16.5)
            };
             \addplot coordinates {
                % arp proxy multiple
                (1000000,20.3)
                (2000000, 12.4)
            };
           \legend{
            A1,A2,A3,A4
          }
        \end{axis}
    \end{tikzpicture}
    \caption{แผนภูมิแสดงผลลัพธ์ประสิทธิภาพด้านการตรวจจับการโจมตีของการทดลองรูปแบบที่ 2}
    \label{img:graph_atk_2}
\end{figure}

จากรูปภาพที่ \ref{img:graph_atk_2} (A1) มีประสิทธิภาพเพิ่มขึ้นเมื่อแพ็กเก็ตเพิ่มขึ้น และมีแนวโน้มประสิทธิภาพเพิ่มขึ้น
(A2) ประสิทธิภาพเพิ่มขึ้นอย่างเห็นได้ชัด และมีแนวโน้มอย่างสูงว่าประสิทธิภาพจะเพิ่มขึ้นอีกหากแพ็กเก็ตเพิ่มขึ้น
(A3) มีประสิทธิภาพเพิ่มขึ้นเมื่อแพ็กเก็ตเพิ่มขึ้น มีแนวโน้มเล็กน้อยว่าประสิทธิภาพจะเพิ่มขึ้นอีก
(A4) มีประสิทธิภาพลดลงเมื่อแพ็กเก็ตเพิ่มขึ้น ซึ่งสวนทางกับเงื่อนไขอื่นอย่างเห็นได้ชัด
\\
% \clearpage
\subsection{สรุปผลการทดลอง}

จากที่ได้กล่าวไปว่า การที่ต้องทดลอง 2 รูปแบบ ก็เพื่อที่จะเปรียบเทียบและเพิ่มความถูกต้อง
เป็นต้นว่า หากทั้งสองรูปแบบสอดคล้องกัน ก็เป็นที่น่าเชื่อว่าผลการทดลองนั้นออกมาถูกต้อง
แต่หากทั้งสองรูปแบบออกมาขัดแย้งกัน ก็เป็นไปได้ว่าการทดลองนี้ไม่ถูกต้อง หรือมีข้อผิดพลาด
ดังนั้น ในการที่จะสามารถบอกได้ชัดเจนเกี่ยวกับผลทดลองนั้น จำเป็นต้องดูทั้ง 2 รูปแบบประกอบกัน

โดยในประสิทธิภาพด้านเวลานั้น จะเห็นได้อย่างชัดเจนว่า (A2) นั้นผลออกมาขัดแย้งกันอย่างชัดเจน จึงกล่าวได้ว่า ผลลัพธ์การทดสอบของ (A2) นั้นมีความน่าเชื่อถือต่ำ
หรืออาจจะมีข้อผิดพลาดในการทดลอง จึงไม่อาจจะนำ (A2) ไปใช้ในการวิเคราะห์อื่นๆ ได้
ส่วนในผลลัพธ์อื่นๆ นั้น จะเห็นว่าทั้ง 2 รูปแบบไม่ได้ขัดแย้งกันอย่างชัดเจน ผลลัพธ์จึงมีความน่าเชื่อถือ
โดย (A4) นั้นมีประสิทธิภาพที่ดีขึ้นเมื่อแพ็กเก็ตมีจำนวนเพิ่มขึ้น ส่วนที่เหลือนั้นไม่ได้มีความต่างอย่างชัดเจน
สามารถกล่าวได้ว่า (A4) ให้ประสิทธิภาพด้านเวลาดีที่สุด

ในประสิทธิภาพด้านการตรวจจับการทดลองนั้น จะเห็นได้อย่างชัดเจนเช่นเดียวกับประสิทธิภาพด้านเวลา ว่า (A2) นั้นมีผลลัพธ์ที่ขัดแย้งกันจากทัง 2 รูปแบบ ดังนั้น (A2) จึงมีความน่าเชื่อถือต่ำ
ในส่วน (A1) นั้น แม้ดูมีแนวโน้มว่าจะขัดแย้งกัน แต่เมื่อพิจารณาผลลัพธ์ในรูปแบบที่ 2 แล้ว ค่ามีการเพิ่มขึ้นเล็กน้อยเพียงเท่านั้น ทำให้ยังไม่ได้มีความขัดแย้งอย่างเห็นได้ชัด
ส่วนอื่นๆ นั้น ถือว่าค่อนข้างมีความสอดคล้องกัน โดยที่ไม่ได้มีข้อใดมีประสิทธิภาพที่ดีกว่าอย่างเห็นได้ชัด หากดูตามกราฟแล้ว พอจะพูดได้ว่า (A3) ให้ประสิทธิภาพโดยรวมที่ดีที่สุด

ประสิทธิภาพด้านการตรวจจับการโจมตีนั้น แม้จะไม่ได้มีความขัดแย้งกันอย่างชัดเจน แต่ผลลัพธ์ก็ไม่ได้เป็นไปในทิศทางเดียวกันมากนัก ดังนั้นจึงต้องมีความระมัดระวังมากเป็นพิเศษในการนำผลลัพธ์ด้านการตรวจจับแพ็กเก็ตไปวิเคราะห์ โดยหากต้องการนำข้อมูลใดๆ ไปวิเคราะห์หรือประกอบการตัดสินใจ แนะนำให้พิจารณาเลือกรูปแบบที่ 2 เป็นอันดับแรก
เพราะมีช่วงจำนวนการจับแพ็กเก็ตมากกว่า
\\
\subsection{ข้อเสนอแนะ}

ในการทดลองนี้ แม้ผู้ทำการทดลองจะออกแบบการทดลองให้มีความรัดกุม และเข้มงวดมากที่สุด
แต่ก็หลีกเลี่ยงไม่ได้ที่จะมีข้อผิดพลาด หรือผลลัพธ์ที่คลาดเคลื่อน แม้แต่ผลลัพธ์ที่ไม่สอดคล้องเองก็ตาม
อย่างไรก็ดี หากต้องการนำผลการทดลองไปใช้ ควรทำการทดสอบและวัดผลด้วยตัวเอง ตามวิธีที่เขียนไว้ หรืออาจจะวิธีการอื่นที่ดีกว่า ซึ่งอาจจะเพิ่มความรัดกุมขึ้นอีก
หรือเพิ่มความแม่นยำของผลลัพธ์ โดยอิงตามทฤษฎีมากขึ้น และควบคุมตัวแปรได้อย่างเรียบร้อย
\end{document}